% !TEX program = lualatex
% Coverage example: a list item containing more than one paragraph. See #20.
%
% A continuation paragraph inside a body item is the only paragraph in a memo that the
% AUTHOR starts and the class never wraps in \noindent. It is therefore the only place a
% stray \parindent could reach the page -- which makes this the regression witness for
% the explicit \parindent=0pt in the class.
%
% It renders flush today, and #20's premise that fourteen \noindent calls "compensate for"
% a non-zero \parindent turned out to be wrong: \RaggedRight (ragged2e) already assigns
% \parindent\RaggedRightParindent, which is zero. The class's explicit setting exists so
% that stops being an accident of another package's side effect. If either is removed,
% these two paragraphs are what moves.
\documentclass{../armymemo}

\address{Organizational Name/Title}
\address{Standardized Street Address}
\address{CITY, STATE~~12345-1234}

\author{John W. Smith}\rank{CPT}\branch{CY}
\officesymbol{ABC-DEF-GH}
\signaturedate{10 April 2019}
\memoline{MEMORANDUM FOR RECORD}
\documentmark{UNCLASSIFIED//FOR OFFICIAL USE ONLY (EXAMPLE ONLY)}
\subject{A memorandum whose items carry more than one paragraph}
\title{COMMANDING}

\begin{document}

\begin{enumerate}
\item This item has two paragraphs. The first one starts the item, so enumitem places it
  and no indent applies.

  This second paragraph is started by the author, not by the class -- the only text in a
  memo that no \verb!\noindent! precedes. It renders flush, and that is the assertion:
  if \verb!\parindent! ever stops being zero, this line moves right by 12pt.
\item A second item, also with two paragraphs.

  Its continuation paragraph is the same assertion a second time, so a single stray
  indent cannot be mistaken for a one-off.
\end{enumerate}

\end{document}
